\subsection{Схема 1: Калышка, пробитая ходовым концом один или более раз}

\graphicspath{{\currentpath}}


\begin{tabular}{
	>{\centering\arraybackslash}m{1cm}
	>{\centering\arraybackslash}m{8cm}
	>{\centering\arraybackslash}m{3cm}
}
1 & \includegraphics[scale=2]{images/simple-0.eps} & простой узел \\
2 & \includegraphics[scale=2]{images/simple-1.eps} & кровавые узлы \\
3 & \includegraphics[scale=2]{images/simple-2.eps} & 
\end{tabular}


\subsection{Схема 2: Обобщение схемы 1 на несколько калышек}

Группа с двумя калышками

\begin{tabular}[b]{
	>{\centering\arraybackslash}m{1cm}
	>{\centering\arraybackslash}m{1cm}
	>{\centering\arraybackslash}m{8cm}
	>{\centering\arraybackslash}m{3cm}
}
2 & 1 & \includegraphics[scale=1.5]{images/double-simple-2-1.eps} & кровавый узел \\
2 & 2 & \includegraphics[scale=1.5]{images/double-simple-2-2.eps} & удвоенный простой \\
2 & 3 & \includegraphics[scale=1.5]{images/double-simple-2-3.eps} & узел с 2+3 шлагами\\
\end{tabular}

Группа с тремя калышками

\begin{tabular}{
	>{\centering\arraybackslash}m{1cm}
	>{\centering\arraybackslash}m{1cm}
	>{\centering\arraybackslash}m{8cm}
	>{\centering\arraybackslash}m{3cm}
}
3 & 1 & \includegraphics[scale=1.5]{images/double-simple-3-1.eps} & кровавый узел \\
3 & 2 & \includegraphics[scale=1.5]{images/double-simple-3-2.eps} & узел с 3+2 шлагами \\
3 & 3 & \includegraphics[scale=1.5]{images/double-simple-3-3.eps} & утроенный простой со шлагом
\end{tabular}